\documentclass[11pt,a4paper]{article}

\usepackage[utf8]{inputenc}
\usepackage[T1]{fontenc}
\usepackage[english]{babel}
\usepackage[margin=2.5cm]{geometry}
\usepackage{graphicx}
\usepackage{booktabs}
\usepackage{listings}
\usepackage{xcolor}
\usepackage{hyperref}
\usepackage{float}
\usepackage{amsmath}
\usepackage{caption}

% Numbers, environment data and tables are generated by analyze.py from results.csv
\input{generated.tex}

\lstset{
  basicstyle=\ttfamily\footnotesize,
  breaklines=true,
  frame=single,
  numbers=left,
  numberstyle=\tiny,
  keywordstyle=\color{blue},
  commentstyle=\color{gray},
  showstringspaces=false
}

\title{Basic Matrix Multiplication in Python, Java and C\\
\large Assignment 1 -- Big Data}
\author{
  Pablo Mart\'in Perdomo\\
  DNI: 78823396G\\
  Group: 17, 42
}
\date{October 16, 2026}

\begin{document}
\maketitle

\begin{abstract}
This report compares the straightforward triple-loop dense matrix multiplication
implemented in Python, Java and C for $n\times n$ integer matrices with
$n$ from 50 up to 1000. Each configuration was measured with five repetitions in
fresh processes, recording median and variability of the runtime, peak resident memory,
and correctness against a reference. All three programs reproduce the
reference results exactly. In the tested environment, runtimes follow the cubic model
approximately (fitted exponents between \ExpJavaCold{} and \ExpJavaWarm{}),
CPython is roughly 90 times slower than C at $n=750$ and exceeds the 60\,s budget at $n=1000$,
while Java is between 1.5 and 1.8 times slower than C at the largest sizes.
Conclusions are limited to the tested implementations, settings, inputs and hardware.
\end{abstract}

\tableofcontents
\newpage

%--------------------------------------------------------------------
\section{Introduction and objective}
Dense matrix multiplication $C = A\cdot B$ for $n\times n$ matrices requires
$n^3$ multiply--add operations with the classical algorithm:
\[
  C_{ij} = \sum_{k=0}^{n-1} A_{ik}\,B_{kj}.
\]
The objective is to compare three implementations of exactly this algorithm
(no optimised library calls) and to investigate the effect of $n$ on runtime
and memory, discussing interpreter, runtime and memory-layout effects only where
the evidence supports them, and comparing the trend with the cubic model.

%--------------------------------------------------------------------
\section{Experimental setup}

\subsection{Hardware and software environment}
\begin{table}[H]
\centering
\begin{tabular}{lp{11cm}}
\toprule
Item & Value \\
\midrule
CPU & \EnvCPU{} (\EnvCores{} virtual cores, cloud VM) \\
RAM & \EnvRAM{} \\
Operating system & \EnvOS{} \\
C compiler & \EnvGcc{} \\
Java / JVM & \EnvJava{} (OpenJDK 64-Bit Server VM, default flags) \\
Python & \EnvPython{} \\
\bottomrule
\end{tabular}
\caption{Execution environment. All measurements were taken on this machine.}
\end{table}
The machine is a shared virtual machine, so some timing noise is expected, especially for
very short runs. NumPy is used only in the test and verification scripts to compute the reference
result; it is never called in any timed code.

\subsection{Implementations}
All three versions use the same algorithm, the same loop order ($i$--$j$--$k$,
accumulating $\sum_k A_{ik}B_{kj}$ in a local variable and storing it once into $C_{ij}$),
the same numeric type family (integers) and the same deterministic inputs
\[
  A_{ij} = (7i+3j+1) \bmod 10 + 1, \qquad B_{ij} = (5i+11j+2) \bmod 10 + 1,
\]
so that the three programs multiply identical matrices and can be compared exactly.
The largest entry of $C$ at $n=1000$ is $10^5$, so 32-bit arithmetic never overflows.

\begin{table}[H]
\centering\small
\begin{tabular}{p{1.5cm}p{4.6cm}p{6.4cm}}
\toprule
Language & Storage layout & Unavoidable differences \\
\midrule
C & \texttt{int**}: array of row pointers, each row a separate \texttt{malloc} block of 32-bit \texttt{int} & Rows contiguous, the whole matrix is not; no bounds checks; native code. \\
Java & \texttt{int[][]}: array of references to \texttt{int[]} rows (32-bit) & Rows contiguous; array bounds checks; code starts interpreted and is JIT-compiled at run time; garbage-collected heap. \\
Python & list of lists; each element is a pointer to an \texttt{int} object & Bytecode interpreter; arbitrary-precision integers; small integers (here the inputs, 1--10) are shared cached objects, but each result element is a separate heap object. \\
\bottomrule
\end{tabular}
\caption{Storage layout and language-specific differences.}
\end{table}

\subsection{Compiler, JVM and interpreter settings}
\begin{itemize}
  \item \textbf{C:} \texttt{gcc -O3 matrix.c -o matrix\_c} (no \texttt{-march=native}, no \texttt{-ffast-math}).
  \item \textbf{Java:} \texttt{javac main.java}; run as \texttt{java main <n> <warmup>} with \emph{default JVM flags}
        (no \texttt{-Xmx}, \texttt{-Xint}, \texttt{-XX} options). Two configurations are reported:
        \emph{Java (cold JVM)}, with 0 warm-up runs, so interpretation and JIT compilation are part of the
        measurement; and \emph{Java (3 warm-up runs)}, where the same multiplication is executed three times
        untimed before the timed run in the same JVM.
  \item \textbf{Python:} CPython \EnvPython{}, pure Python loops (no NumPy, no C extensions); no warm-up is
        applicable.
\end{itemize}
\textbf{Settings held fixed:} algorithm and loop order, input generator, integer type family, matrix sizes, repetitions,
timeout, hardware, compiler flags and JVM flags. \textbf{What differs} between the tested implementations is, by design,
the language implementation (native code, JIT, interpreter).

\subsection{Measurement protocol}
\begin{itemize}
  \item \textbf{Sizes:} $n\in\{50,100,150,200,250,300,400,500,750,1000\}$.
  \item \textbf{Repetitions:} 5 per configuration, each in a \emph{fresh process}; repetitions are interleaved
        across implementations to spread system drift. Total: \NRuns{} successful runs.
  \item \textbf{Time:} only the multiplication is timed, with a monotonic wall clock
        (\texttt{clock\_gettime} in C, \texttt{System.nanoTime()} in Java, \texttt{time.perf\_counter()} in Python).
        Input generation, process start-up and JVM start-up are excluded. We report median, min, max and
        the coefficient of variation (sample std / mean).
  \item \textbf{Memory:} metric = \emph{peak resident set size} (maximum RSS) of the whole process, obtained with the
        small wrapper \texttt{memrun} (\texttt{wait4()} / \texttt{ru\_maxrss}, equivalent to \texttt{/usr/bin/time -f \%M}).
        It includes the JVM and the interpreter themselves.
  \item \textbf{Budget:} 60\,s per run. A run exceeding the budget is killed, recorded as timeout, and larger sizes of
        that configuration are skipped.
  \item \textbf{Correctness during the benchmark:} every run prints a weighted checksum
        $\sum_{i,j}(i+1)(j+1)C_{ij}$ that is compared with a NumPy int64 reference.
\end{itemize}

%--------------------------------------------------------------------
\section{Correctness}
The tests (\texttt{test\_correctness.py}) use exact integer arithmetic, hence the numerical tolerance is \textbf{zero}
(exact equality required).
\begin{enumerate}
  \item \textbf{Hand-computed case} ($n=2$): $A=\begin{pmatrix}2&5\\9&2\end{pmatrix}$,
        $B=\begin{pmatrix}3&4\\8&9\end{pmatrix}$, so $C=\begin{pmatrix}46&53\\43&54\end{pmatrix}$.
  \item \textbf{Full-matrix comparison} with a NumPy reference for $n\in\{1,2,3,5,10,16,33,64\}$
        (including $n=1$ and non-powers of two).
  \item \textbf{Checksum comparison} with the reference for $n=100$ and $n=250$.
\end{enumerate}
\input{tab_correct.tex}
In addition, the checksum of every one of the \NRuns{} benchmark runs matched the reference
(\AllChecksumsOK{} runs correct).

%--------------------------------------------------------------------
\section{Results}

\subsection{Runtime}
\input{tab_median.tex}
\input{tab_var.tex}

Variability is small for long runs (coefficient of variation of a few percent for Python above $n=250$) and large for
very short runs: C below $n=300$ and Java below $n=400$ take milliseconds or less, and the measurements show
coefficients of variation that are often 10--40\% (98\% for the Java warm-up case at $n=100$, caused by a single slow repetition).
Medians are therefore used throughout.

\begin{figure}[H]
\centering
\includegraphics[width=0.95\textwidth]{fig_runtime_linear.png}
\caption{Runtime versus matrix size (linear axes). Python dominates; the other curves are compressed near zero.}
\label{fig:linear}
\end{figure}

\begin{figure}[H]
\centering
\includegraphics[width=0.95\textwidth]{fig_runtime_loglog.png}
\caption{Runtime versus matrix size on log--log axes; error bars show min--max over 5 repetitions. The dotted line has
slope 3 (cubic model) anchored at $n=200$.}
\label{fig:loglog}
\end{figure}

\input{tab_ratio.tex}

\subsection{Growth compared with the cubic model}
\input{tab_exp.tex}
\input{tab_doubling.tex}

\subsection{Memory}
\input{tab_mem.tex}

\begin{figure}[H]
\centering
\includegraphics[width=0.9\textwidth]{fig_memory.png}
\caption{Peak resident set size (median of repetitions) versus matrix size.}
\label{fig:mem}
\end{figure}

\subsection{Limits}
\begin{itemize}
  \item Largest size completed within the 60\,s budget: C \LimitC, Java (cold) \LimitJavaCold, Java (warm-up) \LimitJavaWarm,
        Python \LimitPython.
  \item Failures: \FailList. No memory failures occurred (the largest peak RSS was below 70\,MiB on a machine with \EnvRAM{} of RAM).
  \item The largest size at which all four configurations completed is $n=\NMaxCommon$.
  \item $n=1000$ was the largest size attempted; larger sizes were not tested for C and Java.
\end{itemize}

%--------------------------------------------------------------------
\section{Analysis}

\paragraph{Growth versus the cubic model.}
The algorithm performs $n^3$ multiply--add operations, so doubling $n$ should multiply the runtime by about 8 if the
cost per operation is constant. Over the sizes $n\ge200$ the fitted log--log slopes are
\ExpPython{} (Python), \ExpC{} (C), \ExpJavaCold{} (Java, cold) and \ExpJavaWarm{} (Java, warm-up), all with $R^2>0.998$.
Python and C therefore grow slightly \emph{faster} than $n^3$, with doubling ratios between 7.7 and 10.7 instead of 8.
The trend is cubic to a first approximation, but the constant is not strictly constant.
A plausible explanation is memory-hierarchy pressure: in the $k$ loop, $B$ is read column-wise, i.e. with a stride of one
row between consecutive accesses, and the three matrices exceed the 2\,MiB L2 cache of this machine once $n\approx 500$ or larger
(three $n\times n$ 4-byte matrices occupy about $12\,n^2$ bytes, 3\,MB at $n=500$). This is a hypothesis consistent with
the data, \emph{not} verified: no cache counters were collected and the timer noise of a shared VM is of the same order as
the effect.

\paragraph{Java: start-up and JIT effects.}
The cold-JVM slope (\ExpJavaCold) is below 3 only because small sizes carry a roughly fixed cost
(interpretation and JIT compilation inside the timed region): at $n=50$ the cold time is about $6\times$ the warm-up time,
and the doubling ratio rises from 2.4 to 7.9 as the fixed part becomes negligible. With three warm-up runs the Java time
falls to within about $1.2$--$1.7\times$ of C for $n\ge 200$; the gap between cold and warm Java disappears by
$n=1000$ (1.74\,s versus 1.73\,s, difference within the repetition spread), meaning that warm-up matters for runs shorter than about a second in this setup.
The warm-up slope (\ExpJavaWarm) is the highest observed and is driven by the steps between $n=100$ and $n=500$
(doubling ratios 7.8 to 12.3); I did not isolate its cause (candidates are JIT tier transitions and cache effects),
so no explanation is claimed.

\paragraph{Interpreter effects.}
CPython executes each inner-loop iteration as several bytecodes with dynamic type checks, reference counting and
allocation of a new integer object for each accumulated value, so the cost per multiply--add is much larger than in compiled code.
In the tested setup the median is about 36.9\,s at $n=750$ versus 0.41\,s for C
(about $90\times$) and the ratio Python/C stays between about 85 and 100 for $n\ge100$; the 60\,s budget is exceeded at $n=1000$
(extrapolating the $n=750$ median with the fitted exponent gives roughly 90\,s, not measured).

\paragraph{Memory.}
The measured peak RSS is dominated by the runtime for small $n$ and by the data for large $n$.
C starts at about 1.7\,MiB and reaches 13.2\,MiB at $n=1000$; the three matrices alone need
$3\cdot4\cdot10^6$ bytes $=11.4$\,MiB, which accounts for the growth almost exactly.
The JVM has a baseline of roughly 42\,MiB (heap reservation, class data and JIT) and grows to 57\,MiB (cold)
and 64\,MiB (with warm-up); the extra memory of the warm-up configuration is consistent with the additional result
matrices allocated by the untimed runs, which the garbage collector has not necessarily released.
Python starts at about 9\,MiB and reaches 39\,MiB at $n=750$. This agrees with a rough estimate of 8-byte pointers for $A$, $B$
and $C$ plus one 28-byte integer object per element of $C$ (inputs are small cached integers), i.e.\ $(3\cdot 8 + 28)\,n^2 \approx 29$\,MB $\approx 28$\,MiB at $n=750$,
plus the interpreter baseline. Python's data footprint per index pair $(i,j)$ is therefore about 4 times that of C and Java
(52 versus 12 bytes).

\paragraph{Scope.}
These results apply only to the tested implementations (CPython \EnvPython, OpenJDK with default flags, gcc \,-O3 code),
the specific input generator, the sizes $50\le n\le1000$ and this virtual machine; they do not show that one language is universally
faster. Absolute times on other hardware will differ; the relative ordering and approximate growth rates are the
transferable findings. Language-specific optimisations (loop interchange, blocking, \texttt{-march=native}, multithreading,
\texttt{long}/\texttt{double} types, NumPy) were deliberately excluded and would have to be reported as separate experiments.

%--------------------------------------------------------------------
\section{Threats to validity}
\begin{itemize}
  \item Shared virtual machine with 2 vCPUs: external load may inflate individual runs; the median and min--max are reported for this reason.
  \item Only 5 repetitions per cell; the standard deviation estimates are rough, especially for runs below 10\,ms.
  \item Java and Python results depend on JVM and interpreter versions and defaults (JIT compiler threads run concurrently on the second core).
  \item The 60\,s budget censors the Python series at $n=750$; the exponent fit for Python uses 6 points.
  \item Peak RSS is a process-level metric and cannot separate data from runtime overhead; the decomposition above is an estimate.
\end{itemize}

%--------------------------------------------------------------------
\section{Reproducibility}
The submitted folder contains: \texttt{matrix.c}, \texttt{main.java}, \texttt{matrix.py}, \texttt{memrun.c},
\texttt{benchmark.py}, \texttt{test\_correctness.py}, \texttt{analyze.py}, \texttt{report.tex}, the raw data
\texttt{results.csv}, \texttt{summary.csv}, \texttt{correctness.csv}, \texttt{env.txt} and the benchmark log.
Requirements: \texttt{gcc}, a JDK, Python 3 with \texttt{numpy} and \texttt{matplotlib}, and \texttt{pdflatex}.
\begin{lstlisting}[language=bash,numbers=none]
pip install numpy matplotlib
python3 test_correctness.py          # 33 exact correctness checks
python3 benchmark.py --reps 5 --timeout 60    # writes results.csv, env.txt (about 8 minutes)
python3 analyze.py                   # figures, tables, summary.csv
pdflatex report.tex && pdflatex report.tex
\end{lstlisting}
Re-running on another machine regenerates all tables, figures and environment fields automatically; the numerical
statements in the prose of Section~5 refer to the run included with this report and should be rechecked after a re-run.

\section{Conclusions}
In the tested environment, the three implementations produce identical, exactly verified results, and all grow
approximately cubically in $n$ (fitted exponents 2.9--3.5; doubling ratios between 6.7 and 12.3 for the pairs with $n\ge200$).
At $n=750$, C took 0.41\,s, Java 0.68--0.72\,s and CPython 36.9\,s; CPython did not finish $n=1000$ within 60\,s, whereas C
and Java did in about 1.1\,s and 1.7\,s. Java's cold-start cost is visible below roughly 0.5\,s of work and vanishes at
the largest size. Memory grows quadratically in all three, with very different constants: about 4 bytes per element in C and
Java plus a fixed runtime baseline (2\,MiB for C, 42\,MiB for the JVM), and much more per element in CPython. The slightly
super-cubic growth of C and Python is consistent with cache effects but was not verified.

%--------------------------------------------------------------------
\appendix
\section{Source code}
\subsection{C (\texttt{matrix.c})}
\lstinputlisting[language=C]{matrix.c}
\subsection{Java (\texttt{main.java})}
\lstinputlisting[language=Java]{main.java}
\subsection{Python (\texttt{matrix.py})}
\lstinputlisting[language=Python]{matrix.py}
\subsection{Peak-RSS wrapper (\texttt{memrun.c})}
\lstinputlisting[language=C]{memrun.c}

\end{document}
